\section{Функционалы}

\begin{defn}{}{}
    \( V \) --- векторное пространство над \( F \)

    \( V^* = \{ \xi : V \to F \: | \: \text{\( \xi \) --- линейное} \} = \hom_F (V, F) \) --- двойственное пространство

    \( V^* \) --- тоже векторное пространство

    \( \xi : V \to F \), \( \eta : V \to F \), \( \lambda \in F \)
    \begin{itemize}
        \item
            \( (\xi + \eta) (v) = \xi(v) + \eta(v) \)
        \item
            \( (\lambda \cdot \xi) (v) = \lambda \cdot \xi (v) \)
    \end{itemize}
\end{defn}

\begin{example}{}{}
    \begin{enumerate}
        \item
            \( (\RR^n)^* \) --- пространство строк
        \item
            \( \RR[x] \to \RR \), \( \lambda \in \RR \)

            \( f \mapsto f(\lambda) \)
        \item
            \( C[a, b] \to \RR \)

            \( f \mapsto \int\limits_{a}^{b} f(x) dx \)
        \item
            \( M_n (\RR) \to \RR \)

            \( A \mapsto \tr (A) \)
    \end{enumerate}
\end{example}

\begin{lemma}{}{}
    \( V \) --- векторное пространство над \( F \)

    \( e_1, \ldots, e_n \) --- базис \( V \)

    Тогда:
    \begin{enumerate}
        \item
            \( \exists! \xi_1, \ldots, \xi_n \in V* \)

            \(
                \xi_i (e_j) = \begin{cases}
                    1 \quad i = j 
                    \\
                    0 \quad i \neq j
                \end{cases}
            \)
        \item
            \( \xi_1, \ldots, \xi_n \) --- базис
        \item
            \( \xi \in V^* \ \xi = \xi (e_1) \xi_1 + \ldots + \xi (e_n) \xi_n \)

            \( v \in V \)

            \( v = \xi_1 (v) e_1 + \ldots + \xi_n(v) e_n \)
    \end{enumerate}
\end{lemma}

\begin{defn}{}{}
    \( (\xi_1, \ldots, \xi_n) \) --- двойственный базис к \( (e_1, \ldots, e_n) \).
\end{defn}

\begin{enumerate}
    \item
        Мы знаем, что базисный набор можно отправить куда угодно,
        причем будет существовать единственное линейное отображение с таким свойством.
        Отсюда сразу вытекает существование и единственность \( \xi_i \)
    \item
        Проверим, что наш набор является ЛНЗ порождающим.

        \begin{itemize}
            \item
                ЛНЗ:

                \( \lambda_1 \xi_1 + \ldots + \lambda_n \xi_n = 0 \) (в \( V^* \))

                \( \forall v \in V \):

                \( (\lambda_1 \xi_1 + \ldots + \lambda_n \xi_n)(v) = 0(v) \) (в \( F \))

                Преобразуем обе части:

                \( \lambda_1 \xi_1 (v) + \ldots + \lambda_n \xi_n(v) = 0 \)

                Подставим \( e_i \), получим:

                \( \lambda_i \cdot 1 = 0 \Rightarrow \lambda_i = 0 \)
            \item
                Порождающий:

                Возьмем \( \xi \in V^* \)

                Покажем, что \( \xi = \xi (e_1) \cdot \xi_1 + \ldots + \xi (e_n) \cdot \xi_n \) (в \( V^* \))

                \( \xi(v) = (\xi (e_1) \cdot \xi_1 + \ldots + \xi (e_n) \cdot \xi_n)(v) \)

                \( \xi(v) = \xi (e_1) \cdot \xi_1 (v) + \ldots + \xi (e_n) \cdot \xi_n (v) \)

                Поскольку у нас линейные функции, достаточно проверить только на базисных векторах:

                \( \xi(e_i) = \xi (e_1) \cdot \xi_1 (e_i) + \ldots + \xi (e_n) \cdot \xi_n (e_i) \)

                \( \xi(e_i) = \xi (e_i) \cdot 1 \) --- ура, сошлось
        \end{itemize}
    \item
        Проверим в лоб:

        \( v = \lambda_1 e_1 + \ldots + \lambda_n e_n \)

        \( \xi_i (v) = \xi_i(\lambda_1 e_1 + \ldots + \lambda_n e_n) \)

        \( \xi_i (v) = \lambda_1 \xi_i (e_1) + \ldots + \lambda_n \xi_ i(e_n) \)

        \( \xi_i (v) = \lambda_i \cdot 1 \) --- ура, доказали, что \( \xi_i (v) = \lambda_i \)
\end{enumerate}

\begin{coroll}{}{}
    Если \( \dim V = n < \infty \),
    то \( \dim V^* = n < \infty \).

    Если же \( \dim V = \infty \),
    то если вы верите в аксиому выбора, то \( \dim V^* > \dim V \),
    а если не верите, то данный результат нельзя ни доказать, ни опровергнуть.
\end{coroll}
